\documentclass[11pt]{article}
% Shared preamble for the hanzo-kernel Paper 07 series.
% One style, one vocabulary, inputted by every paper so the series reads as one voice.
\usepackage[margin=1in]{geometry}
\usepackage{booktabs}
\usepackage{amsmath}
\usepackage{amssymb}
\usepackage{graphicx}
\usepackage{xcolor}
\usepackage{hyperref}
\hypersetup{colorlinks=true, linkcolor=blue, urlcolor=blue, citecolor=blue}
\usepackage{enumitem}
\usepackage{listings}
\usepackage{array}

\definecolor{kw}{rgb}{0.13,0.29,0.53}
\definecolor{cmt}{rgb}{0.40,0.40,0.40}
\definecolor{str}{rgb}{0.16,0.50,0.16}
\lstset{
  basicstyle=\ttfamily\footnotesize,
  keywordstyle=\color{kw}\bfseries,
  commentstyle=\color{cmt}\itshape,
  stringstyle=\color{str},
  showstringspaces=false,
  breaklines=true,
  columns=fullflexible,
  frame=single,
  framesep=4pt,
}
\lstdefinelanguage{Rust}{
  morekeywords={fn,let,mut,pub,struct,impl,trait,enum,match,if,else,for,while,loop,return,use,mod,crate,self,super,async,await,where,type,const,static,unsafe,ref,move,dyn,as,true,false},
  morecomment=[l]{//}, morecomment=[s]{/*}{*/}, morestring=[b]", sensitive=true,
}

\newcommand{\x}{\ensuremath{\times}}
\newcommand{\dpa}{\texttt{dp4a}}
\newcommand{\hk}{\texttt{hanzo-kernel}}
\newcommand{\code}[1]{\texttt{#1}}

% Absorb minor prose overfulls without rivers; keep line-breaking sane.
\tolerance=800
\emergencystretch=3em

\usepackage{amsthm}

\newtheorem{proposition}{Proposition}

% Honesty labels: distinguish what is built from what is designed.
\newcommand{\built}{\textcolor{str}{\small\textbf{[built]}}}
\newcommand{\partialb}{\textcolor{kw}{\small\textbf{[partial]}}}
\newcommand{\design}{\textcolor{cmt}{\small\textbf{[design]}}}

\title{\textbf{The Three-Plane Compute Fabric:\\
Consensus-Managed, ZAP-Transported\\ Decentralized AI Inference}}
\author{Hanzo AI --- ML Systems}
\date{Architecture Blueprint --- \emph{One Source, Every Backend} series, 2026}

\begin{document}
\maketitle

\begin{abstract}
A heterogeneous GPU fleet already runs a single model across ROCm, CUDA, and Metal
today: the Hanzo engine's socket-ring shards a model tensor- and pipeline-parallel across
mismatched accelerators over a canonical \code{f32} wire, and a \code{hanzo cluster}
command onboards a worker by scanning a QR token. But that ring is configured by a static
JSON file handed out of band, and its bytes cross the wire as \emph{plaintext TCP} between
nodes that never authenticate each other. It is a LAN appliance, not a network. This paper
is the blueprint for the missing two-thirds: turning that ring into a permissionless,
byzantine-fault-tolerant AI inference \emph{network} that any developer --- or any
public-cloud GPU miner --- can join and be paid for useful work. The thesis is one of
\emph{decomplection}: a serving cluster has three concerns that monolithic stacks braid
together and we hold apart as three orthogonal, independently-complete planes. The
\textbf{control plane} is a Lux \emph{Quasar} consensus chain: membership, permissionless
miner admission, layer placement, and config distribution become BFT-agreed state, with the
Photon committee-selection primitive (Fisher--Yates over a luminance-weighted population)
reused to sample GPU nodes and luminance re-bound from consensus-vote success to
\emph{validated inference contribution}. The \textbf{transport} is \emph{ZAP}, the Lux
zero-copy wire, whose post-quantum handshake (\code{WrapPQ} over an authenticated
\code{Session} exposing a verified 32-byte \code{PeerID}) carries \emph{both} consensus
control messages \emph{and} tensor activations over one PQ-authenticated channel --- retiring
the plaintext ring sockets and converting the QR join token from a bearer secret into a
handshake credential. The \textbf{data plane} is the existing inference ring, unchanged in
kind: it receives its topology from consensus instead of a JSON file and moves its tensors
over ZAP instead of raw TCP. We are scrupulous about the built/designed line. Built and
cited to source: the cross-vendor ring, the \code{hanzo cluster} scan-to-join CLI, the ZAP
Go wire with its PQ handshake, the Quasar engine with Photon/Wave/Focus and BLS$+$Pulsar$+$ML-DSA
threshold signing, and the ZAP measurements ($17\times$ serialization and $345\times$ parse
versus JSON-RPC, $\sim$$4.5\times$ versus gRPC with zero marshal/unmarshal). Designed, and
labelled as such: consensus-driven membership, the ZAP-swapped ring transport, miner
admission, and result verification by threshold-signing over inference outputs. No benchmark
here is invented; every measured number is attributed to the code that produced it. The
contribution is the architecture and its honest status, offered as a blueprint the wider
AI-infrastructure community can build on.
\end{abstract}

\section{Introduction}
\label{sec:intro}

The dominant shape of a model-serving stack is a monolith: one operator, one datacenter, one
process tree that owns scheduling, transport, and tensor math all at once. vLLM, TGI, and
Triton each fuse ``who is in the cluster,'' ``how bytes move between ranks,'' and ``what the
GPUs compute'' into a single trusted binary on a private network. That fusion is fine when the
operator owns every GPU. It fails the moment the fleet is \emph{heterogeneous} (a ROCm mini-PC,
a rented CUDA instance, an Apple laptop), \emph{permissionless} (a stranger's GPU wants to
contribute and be paid), or \emph{untrusted} (the node holding your activations is not yours).
None of those three failures is a math problem; each is a \emph{concerns} problem. The cluster
braids together things that want to be separate.

This paper takes the opposite stance, the one our engineering culture states as a law: one and
only one way to do everything, composable and orthogonal. A decentralized inference cluster has
exactly three concerns, and we give each its own plane:
\begin{itemize}[leftmargin=1.4em,itemsep=2pt]
\item \textbf{Control} --- who is in the cluster, who may join, which node holds which layers,
and what the run configuration is. This is agreement among mutually distrusting parties: a
\emph{consensus} problem.
\item \textbf{Transport} --- how a byte gets from one node to another, and how each end knows
who the other is. This is an \emph{authenticated wire} problem.
\item \textbf{Data} --- the actual tensor exchange: activations forward, gradients or
all-reduce sideways. This is a \emph{numerics and topology} problem.
\end{itemize}
The claim is that these three are genuinely orthogonal, that Hanzo and Lux already ship a
best-in-class primitive for each, and that composing them --- rather than reinventing a
monolith --- yields something no monolith can be: a GPU network anyone can join. The control
primitive is Lux \emph{Quasar} consensus~\cite{quasar,luxconsensus}; the transport is
\emph{ZAP}~\cite{zap,lp201}; the data plane is the Hanzo cross-backend inference
engine~\cite{infstack}.

\paragraph{The honesty frame.} This is an architecture paper, and its credibility rests on not
blurring what runs today with what is designed. We adopt three labels, used inline throughout
and collected in the status table of \S\ref{sec:status}: \built{} for code on \code{main},
cited to its file; \partialb{} for a primitive that is compiled and exercised but not yet wired
end-to-end; and \design{} for a proposed composition. The measured numbers we quote (ring
parity across vendors, ZAP throughput and latency, Photon and luminance costs) are all \built{},
each attributed to the benchmark that produced it. The network-level assembly --- consensus as
the live scheduler, the ring running over ZAP, miner admission and payment --- is \design{},
and never dressed as more.

\paragraph{Contribution.} We (i) inventory the built substrate with source citations
(\S\ref{sec:today}); (ii) state the three-plane architecture and prove its concerns are
separable (\S\ref{sec:planes}); (iii) give the threat model and a result-integrity scheme for
byzantine miners (\S\ref{sec:threat}); (iv) bind the economic layer to the existing luminance
primitive (\S\ref{sec:econ}); (v) lay out a staged migration from the static-config LAN cluster
to a public network (\S\ref{sec:migration}); and (vi) publish an honest built-versus-designed
status table (\S\ref{sec:status}) and the limitations (\S\ref{sec:limits}).

\section{The substrate that exists today}
\label{sec:today}

Before proposing a network we account for what is already on \code{main}. Every claim in this
section is \built{} and cites the file that implements it. The point of the inventory is that
the three planes are not vaporware waiting to be written: each has a mature primitive, and the
work is composition, not invention.

\subsection{Data plane: the heterogeneous inference ring}
\label{sec:ring}

The Hanzo engine already runs one model split across several GPUs of \emph{different vendors}.
Its \code{distributed} module (\code{hanzo-engine/src/distributed.rs},
\code{hanzo-quant/src/distributed/}) provides two backends behind one \code{Comm} abstraction:
NCCL for homogeneous NVIDIA, and a portable \emph{ring} backend for everything else. The ring
does tensor-parallel and pipeline-parallel sharding by exchanging activations with its
right-hand neighbor, rank $r$ dialing rank $(r{+}1) \bmod N$ in an all-reduce cycle
(\code{hanzo-quant/src/distributed/socket.rs}). Cross-vendor generation works because the ring
wire is a \emph{canonical} \code{f32} format and reductions take a backend-independent
CPU-\code{f32} path with a canonical tie-break (\code{hanzo-engine/src/ops.rs}), so a ROCm box
and a CUDA box agree bit-for-bit on the reduced tensor rather than each rank reinterpreting the
other's raw 16-bit words in the wrong \code{f16}-versus-\code{bf16} layout. The topology is a static struct read from a JSON file
named by the \code{RING\_CONFIG} environment variable (Listing~\ref{lst:ringcfg}): master
address, per-rank ports, this rank's index, and the world size. The cross-backend engine
underneath --- five backends, the full GGUF quant zoo, decode parity with \code{llama.cpp} --- is
the subject of the companion stack paper~\cite{infstack}; here it is simply the data plane.

\begin{lstlisting}[caption={The data-plane topology today: a static \code{RING\_CONFIG} JSON.
Every node needs the full member list up front; nothing authenticates a peer.},label={lst:ringcfg}]
{ "master_ip": "10.0.0.1", "master_port": 8000,
  "port": 8001, "right_ip": "10.0.0.2", "right_port": 8001,
  "rank": 0, "world_size": 4 }
\end{lstlisting}

Two properties of this substrate matter for everything below. First, it is \emph{correct across
vendors} --- the hard part, cross-backend numeric agreement, is solved. Second, it is
\emph{unauthenticated plaintext}: the ring exchanges raw bytes over \code{std::net::TcpStream}
with \code{write\_all}/\code{read\_exact} and ships its control messages as \code{serde\_json}
over the same socket. It is safe only on a trusted private mesh. That is the gap the transport
plane closes.

\subsection{The onboarding UX: \code{hanzo cluster} scan-to-join}
\label{sec:clustercli}

Handing every node a JSON file out of band does not scale past a lab. The \code{hanzo cluster}
subcommand (\code{hanzo-cli/src/cluster.rs}) already replaces it with a scan-to-join flow.
\code{hanzo cluster init} on the head (rank 0) enumerates the ring, prints a \emph{join token}
and its QR code, and launches; \code{hanzo cluster join} on a worker decodes the token and
launches as a ring daemon. The token is a \code{bs58}-encoded JSON \code{JoinToken}
(Listing~\ref{lst:token}) --- a nonce, the head address, the world size, the ordered ring
addresses, and the model --- compact enough to fit in one QR or one paste. This is the correct
\emph{shape} of onboarding: a worker needs one string to join. But observe what the token is
today: a \emph{bearer credential in cleartext}. Anyone who photographs the QR can join the ring,
and the ring cannot tell an invited worker from an intruder, because the token carries topology,
not identity. Holding that observation is the whole of \S\ref{sec:threat}.

\begin{lstlisting}[language=Rust,caption={The join token (\code{hanzo-cli/src/cluster.rs}):
\code{bs58(JSON)} so it fits one QR. It is topology, not identity --- a bearer secret today.},label={lst:token}]
struct JoinToken {
    nonce: u64,           // ring id, printed {:016x}
    master: String,       // rank-0 (head) reachable IP
    world_size: usize,
    addrs: Option<Vec<String>>, // full ordered ring [rank0..rankN-1]
    model: String,
}
\end{lstlisting}

\subsection{Transport: ZAP, the zero-copy post-quantum wire}
\label{sec:zap}

ZAP (Zero-copy Application Protocol, \code{github.com/luxfi/zap}) is the Lux universal
wire~\cite{zap,lp201}. Its serialization is zero-copy: a message is read directly out of the
network buffer with no parse and no allocation (measured \built{} parse $2.9$~ns/op at zero
allocations; build $44$~ns/op). Against the protocols an AI stack would otherwise reach for,
the wire eliminates the marshal/unmarshal step entirely: $17\times$ faster serialization and
$345\times$ faster parse than JSON-RPC (the MCP transport), $\sim$$157\times$ faster than
protobuf on deserialization ($21$~ns versus $3231$~ns, zero allocations), and $\sim$$4.5\times$
faster than gRPC with the whole codec layer deleted rather than optimized~\cite{zap,lp201,luxconsensus}.
End-to-end it sustains $114$K transactions/s on a single connection and $376$K across $20$
parallel connections~\cite{luxconsensus}.

For a public network the decisive feature is not speed but \emph{authenticated identity}. ZAP's
PQ layer (\code{conn\_pq.go}) wraps an established connection \emph{after} a post-quantum
handshake: \code{handshake.Initiator.Run(tcp)} produces a \code{Session}, and \code{WrapPQ(tcp,
sess)} returns an ordinary \code{net.Conn} that AEAD-frames every write, zeroes its keys on
close, and --- critically --- exposes \code{PeerID() [32]byte}, ``the verified peer identity for
callers that need to route on it (e.g.\ validator-set membership checks).'' That one method is
the hinge of the entire security argument (\S\ref{sec:threat}): the transport already hands the
control plane a cryptographically verified identity for the node on the other end.

Inside the Hanzo engine, ZAP is a first-class module (\code{hanzo-engine/src/zap/}). Its
message and adapter surface is compiled and exercised in-process --- \code{ZapInferenceServer}
translates ZAP \code{ChatRequest}s into the engine's real \code{Request} channel and streams
\code{ChatChunk}s back, satisfying the same \code{ZapInference} trait a networked client would
(\code{server.rs}, \code{client.rs}). We mark this \partialb{}, honestly: the out-of-process
\code{zap://} transport is a documented stub that carries no \code{capnp} build dependency yet
(``flipping on the wire is mechanical,'' \code{transport.rs}), and a few methods
(\code{tokenize}, \code{detokenize}, \code{re\_isq}) return ``not yet wired.'' The wire exists
in Go and is production-deployed across Lux consensus, MPC, and DEX systems; the engine-side
Rust surface is built but not yet turned on over a socket.

\subsection{Control: the Quasar consensus engine}
\label{sec:quasar}

Lux \emph{Quasar} (\code{github.com/luxfi/consensus}) is a DAG-native, post-quantum BFT
consensus engine~\cite{quasar,luxconsensus}. It is built from small, orthogonal sub-protocols,
four of which this paper reuses directly:
\begin{itemize}[leftmargin=1.4em,itemsep=2pt]
\item \textbf{Photon} (\code{protocol/photon}) selects a $K$-of-$N$ committee each round by
Fisher--Yates shuffle over the node population, drawing with a rejection-sampled uniform integer
(\code{emitter.go}) to remove the modulo bias that a grinding adversary could exploit on
committee sampling. The population is weighted by \emph{luminance}.
\item \textbf{Luminance} (\code{photon/luminance.go}) is a per-node reputation in ``lux'' units:
each node starts at $100$, a successful vote multiplies its brightness by $1.1$ (capped at
$1000$), a failure by $0.9$ (floored at $10$), and \code{Brightness} normalizes to a selection
weight in $[0.1, 10]$. It is a cheap, bounded, exponential reputation --- a brightness update
costs $72$~ns \built{}.
\item \textbf{Wave} (\code{protocol/wave}) runs per-round threshold voting with Fast
Probabilistic Consensus; a vote round costs $3.38$~\textmu s and a $K$-of-$N$ Photon selection
$3.03$~\textmu s \built{}.
\item \textbf{Quasar} (\code{protocol/quasar}) seals decisions with threshold signatures across
three independent cryptographic paths --- BLS12-381 (48-byte aggregate), Pulsar Ring-LWE
threshold, and ML-DSA-65 (FIPS 204, rolled to a 192-byte Groth16 proof) --- each independently
toggleable, all three in parallel under \code{TripleSignRound1}. A QuasarCert verifies in
$\sim$$2$--$5$~ms on CPU, GPU-batchable $10$--$100\times$~\cite{luxconsensus}.
\end{itemize}
Quasar's own inter-node transport is already ZAP, not gRPC or libp2p (\code{consensus/LLM.md}).
Control and transport are, in the Lux stack, \emph{already} composed; what is missing is binding
that composition to the \emph{GPU} population instead of the validator population.

\section{The three-plane architecture}
\label{sec:planes}

The blueprint is a single sentence: \emph{run the inference ring's topology as Quasar consensus
state, and move every byte --- control and tensor alike --- over ZAP.} Figure~\ref{fig:planes}
is the whole system. What follows gives each plane its interface and states, as a proposition,
why the three do not leak into each other.

\begin{figure}[t]
\centering
\small
\begin{tabular}{@{}l@{}}
\toprule
\textbf{Control plane} --- Quasar consensus (BFT, post-quantum) \hfill \design{} \\
\quad Photon committee over GPU nodes $\cdot$ luminance = useful-work weight \\
\quad membership $\cdot$ miner admission $\cdot$ layer placement $\cdot$ config $\cdot$ result attest \\
\midrule
\textbf{Transport plane} --- ZAP PQ wire \hfill \built{}\,/\,\partialb{} \\
\quad \code{WrapPQ} session $\cdot$ verified \code{PeerID} $\cdot$ AEAD frames \\
\quad one channel for BOTH control messages AND tensor activations \\
\midrule
\textbf{Data plane} --- Hanzo inference ring \hfill \built{} \\
\quad tensor/pipeline parallel $\cdot$ canonical \code{f32} all-reduce $\cdot$ 5 backends \\
\quad topology handed down by control; bytes carried by transport \\
\bottomrule
\end{tabular}
\caption{Three orthogonal planes. Each is independently complete and separately citable to
code; the network is their composition, not a new monolith. Labels mark built versus designed
per \S\ref{sec:status}.}
\label{fig:planes}
\end{figure}

\subsection{Control plane: consensus as the scheduler}
\label{sec:control}

The control plane replaces the static \code{RING\_CONFIG} JSON (Listing~\ref{lst:ringcfg}) and
the out-of-band QR token with \emph{agreed state on a Quasar chain}. The chain's VM state
\emph{is} the cluster's control surface \design{}:
\begin{itemize}[leftmargin=1.4em,itemsep=2pt]
\item \textbf{Membership} --- the set of GPU nodes, each identified by its ZAP \code{PeerID}
(an ML-DSA public key), with advertised capability (backend, VRAM, resident models, measured
decode/prefill throughput). Joining and leaving are transactions, so the member list is
BFT-agreed rather than hand-distributed.
\item \textbf{Miner admission} --- a permissionless join. A public-cloud GPU submits a join
transaction carrying its \code{PeerID} and a stake bond; admission is a consensus decision, not
an operator's SSH session. This is the step that turns ``my four boxes'' into ``anyone's GPU''
(\S\ref{sec:econ}).
\item \textbf{Placement / routing} --- which node holds which contiguous layer range of which
model. Pipeline-parallel assignment (node $A$ holds layers $0$--$7$, $B$ holds $8$--$15$) and
the ring order become an on-chain placement map. A request is routed by reading that map, not by
trusting a JSON file.
\item \textbf{Config distribution} --- model id, quantization, sampling defaults, ring
\code{world\_size}: the fields of \code{JoinToken} and \code{RingConfig} become chain state, so
every node reads an identical, agreed configuration.
\end{itemize}

The mechanism is not new code so much as a \emph{re-population} of Photon. Photon selects a
committee by weighted Fisher--Yates over a node set; point it at the \emph{GPU} population and it
selects a committee of GPU nodes --- to ratify a placement, to admit a miner, or to verify a
result (\S\ref{sec:threat}). Luminance, today a function of consensus-vote success, is re-bound
to \emph{validated inference contribution}: a node whose outputs pass verification brightens and
is preferentially selected and rewarded; a node that stalls or returns wrong tensors dims and is
sampled out. The primitive is unchanged --- an exponential, bounded, cheap reputation --- only
the event that drives \code{Illuminate(id, success)} moves from ``voted with the quorum'' to
``did correct inference work.'' Crucially, \emph{writes to consensus are rare and
slow-changing}: placement and membership shift on the order of nodes joining and models loading,
not on the order of tokens generated. The hot path --- a forward pass through the ring --- never
touches consensus.

\subsection{Transport plane: one PQ wire for control and tensors}
\label{sec:transport}

Every inter-node byte moves over ZAP. This is deliberately \emph{one} wire for two kinds of
traffic that a monolith would split across two stacks:
\begin{itemize}[leftmargin=1.4em,itemsep=2pt]
\item \textbf{Control messages} --- consensus votes, placement updates, admission requests ---
which Quasar already speaks over ZAP.
\item \textbf{Tensor data} --- the activation exchange and all-reduce that the ring performs
every micro-step --- which today crosses raw TCP and, under this design, crosses the same
authenticated \code{WrapPQ} \code{net.Conn} instead.
\end{itemize}
Unifying them on one wire is the payoff of decomplection: the ring does not need to know it is on
a network of strangers, only that it holds a \code{net.Conn}; the transport does not need to know
whether the bytes are a vote or an activation. ZAP's zero-copy framing is what makes carrying
tensors over an authenticated session affordable --- there is no marshal/unmarshal tax on the
activation exchange, only an AEAD seal over a buffer the ring already owns (the $64$~KiB
\code{MaxPQRecord} chunking of \code{conn\_pq.go} maps naturally onto activation-sized writes).
The identity that the handshake establishes, \code{PeerID}, is the same value the control plane
keys membership on: the node you complete a ZAP handshake with \emph{is} the node consensus
admitted, by construction, with no separate PKI to reconcile.

\subsection{Data plane: unchanged in kind}
\label{sec:data}

The data plane is the least-changed and that is the point. The ring's numerics --- canonical
\code{f32} reduction, tensor/pipeline sharding, the five-backend kernels --- are correct today
(\S\ref{sec:ring}) and stay exactly as they are. Two wires change underneath them: the topology
arrives from the control plane's placement map rather than \code{RingConfig::load()}, and the
\code{TcpStream} is a ZAP \code{net.Conn} rather than a bare socket. The ring code that computes
does not move. This is composition over rewrite: we are not building a decentralized inference
engine, we are giving the inference engine we have a control plane and an authenticated
transport.

\subsection{Why three planes, stated precisely}
\label{sec:whythree}

\begin{proposition}[Orthogonality of the three planes]
\label{prop:ortho}
The control, transport, and data planes can each be replaced or scaled independently without
modifying the other two, because they share only two narrow, typed interfaces: the control plane
emits a \emph{placement map} keyed by \code{PeerID}, and the transport plane emits an
authenticated \code{net.Conn} whose \code{PeerID} matches. The data plane consumes both and
exposes neither consensus nor cryptography; the transport consumes an identity and exposes a
byte stream, knowing nothing of tensors; the control consumes admission and contribution events
and exposes agreed state, knowing nothing of wire framing.
\end{proposition}

\begin{proof}[Argument]
By interface inspection. The data plane's dependencies are exactly (i) a topology --- rank,
world size, neighbor address, layer range --- and (ii) a \code{net.Conn} per neighbor. Neither
names a consensus type or a cipher. Swap Quasar for any state machine that produces the same
placement map and the ring is unaffected; swap ZAP for any authenticated \code{net.Conn} with a
stable peer identity and the ring is unaffected; swap the ring's kernels (add a backend, change
the quant) and neither consensus nor transport observes it. The only cross-plane invariant is
that the \code{PeerID} the transport authenticates equals the \code{PeerID} the control plane
admitted --- a single equality, checked once at connection setup, which is precisely what
\code{WrapPQ}'s \code{PeerID()} was built to support (``validator-set membership checks,''
\code{conn\_pq.go}). Three concerns, two interfaces, one shared value. \qed
\end{proof}

The contrast with a monolithic serving stack is now sharp. In vLLM or Triton, ``who is in the
cluster'' is a launch-time flag baked into the process, ``how bytes move'' is a hardcoded NCCL
or gRPC channel with no peer authentication, and ``what computes'' is the same binary --- the
three concerns are one artifact, so you cannot add a permissionless member, cannot authenticate
a stranger's GPU, and cannot verify a result without forking the whole thing. Our planes are
three artifacts with two seams. That is the difference between a cluster and a network.

\section{Threat model and result integrity}
\label{sec:threat}

A permissionless network inherits adversaries a private cluster never has. We state the model,
the gaps in today's substrate, and the design that closes them.

\paragraph{Assets and adversaries.} The assets are: the integrity of an inference result (the
tokens a user is billed for must be the model's real output), the confidentiality and
authenticity of activations in flight, the correctness of the placement map, and the fairness of
reward. The adversaries are: a network attacker on the path between nodes (public-cloud tenants
share fabric); an \emph{intruder} who obtained a join token; and a \emph{byzantine miner} who
joined honestly but returns wrong, fabricated, or lazily-cached tensors to earn reward without
doing the work.

\paragraph{Gap 1: the wire is plaintext (\built{} problem).} Today the ring is raw TCP
(\S\ref{sec:ring}). A path attacker can read activations, inject a forged tensor into the
all-reduce, or man-in-the-middle a neighbor. \textbf{Closure \design{}:} the ZAP PQ handshake.
After \code{WrapPQ}, every frame is AEAD-sealed and every peer is a verified ML-DSA identity;
injection and eavesdropping both fail against a post-quantum authenticated channel. The tensor
plane inherits the same wire security consensus already relies on.

\paragraph{Gap 2: the join token is a bearer secret (\built{} problem).} The QR token
(Listing~\ref{lst:token}) is cleartext topology; whoever holds it joins (\S\ref{sec:clustercli}).
\textbf{Closure \design{}:} make the token a \emph{handshake credential}, not the thing that
grants access. Under the design the token carries the ring's \code{PeerID} set (or a commitment
to it); a worker presenting the token must then \emph{complete a ZAP PQ handshake} proving
possession of an admitted ML-DSA key. A photographed QR no longer joins anyone --- it only names
who is allowed to, and access is gated on the handshake the intruder cannot forge. The token
degrades from a secret to a public address book, which is the correct security posture: identity
is the credential, topology is not.

\paragraph{Gap 3: a miner may lie about its work (the hard one).} Authentication proves
\emph{who} sent a tensor, not that the tensor is \emph{correct}. A byzantine miner holding layers
$8$--$15$ can return plausible-looking garbage and collect reward. This is the genuinely new
problem, and consensus is the tool. We propose three composable mechanisms, in increasing cost,
scheduled by the control plane \design{}:
\begin{enumerate}[leftmargin=1.5em,itemsep=2pt]
\item \textbf{Threshold attestation of outputs.} A completed inference (or a placement-epoch's
batch of them) is sealed with a QuasarCert: a Photon committee of GPU nodes threshold-signs the
$(\text{request}, \text{output})$ pair with BLS$+$Pulsar$+$ML-DSA (\S\ref{sec:quasar}). A billed
result carries a certificate that a supermajority attested to it, giving the user cryptographic
non-repudiation and making a single lying miner insufficient to certify a bad output.
\item \textbf{Redundant re-execution / spot checks.} The scheduler occasionally assigns the same
layer range to two independent miners and compares --- cheap because the ring wire is already
\emph{canonical} \code{f32} (\S\ref{sec:ring}), so two honest miners on different vendors produce
\emph{bit-identical} tensors and any disagreement is a proof of fault, not a rounding artifact.
The canonical-\code{f32} property, built for cross-vendor correctness, doubles as the substrate
for verifiable redundancy. A miner caught disagreeing with the honest majority is dimmed by
luminance and slashed.
\item \textbf{Reputation as the economic backstop.} Luminance (\S\ref{sec:quasar}) makes
correct work compound and incorrect work decay: a node that fails verification multiplies its
brightness by $0.9$ and is sampled into fewer committees and paid for less work, converging its
influence and income toward zero. Reputation is the slow, cheap layer under the fast, expensive
cryptographic ones.
\end{enumerate}
Byzantine fault tolerance is a property the control plane already has (Quasar tolerates $<1/3$
adversarial stake); the contribution here is \emph{extending it from block finality to inference
finality} by threshold-signing over outputs and exploiting canonical \code{f32} for cheap
re-execution. We are explicit that this is \design{}: the threshold-signing primitive is built,
its application to inference outputs is proposed.

\paragraph{Non-goal: confidentiality from the operator.} This design authenticates and verifies;
it does not, by itself, hide activations from the miner that computes them. Workloads needing
that require a hardware TEE tier (H100/H200 confidential computing, TDX, SEV-SNP), an orthogonal
extension developed in the tenant-cloud companion~\cite{tenantcloud} and out of scope here.

\section{The economic layer: luminance as useful-work weighting}
\label{sec:econ}

A permissionless network needs an incentive, and the incentive must be \emph{useful-work},
because paying for wasted computation (proof-of-work hashing) is exactly what an inference
network should not do. The luminance primitive is already a useful-work meter in miniature: it
rewards nodes for correct participation and decays nodes for failure. The economic design
\design{} binds three quantities to it:
\begin{itemize}[leftmargin=1.4em,itemsep=2pt]
\item \textbf{Admission by stake.} A miner joins by bonding stake in a transaction
(\S\ref{sec:control}). Stake is the sybil cost --- spinning up a thousand fake GPUs is a
thousand bonds --- and the slashing target when verification (\S\ref{sec:threat}) catches a
fault.
\item \textbf{Reward by validated contribution.} Payment for an inference accrues to the miners
whose outputs were attested, weighted by luminance and by measured work (tokens $\times$ model
size $\times$ layer share). Because luminance rises with correct work and falls with incorrect
work, expected income tracks \emph{honest throughput}, and a lazy or lying miner earns less than
its electricity costs.
\item \textbf{Routing by brightness.} The Photon committee that ratifies placement and verifies
results draws with luminance weight, so brighter (more reliably correct) nodes are preferentially
given work --- reputation and revenue reinforce, which is what makes ``contribute a GPU, get
paid'' a stable equilibrium rather than a race to the bottom.
\end{itemize}
The mechanism reuses one existing primitive (luminance) for reputation, one existing primitive
(stake-weighted BFT) for sybil-resistance and slashing, and one existing primitive (Photon) for
work routing. The realized system is ``any dev, any miner, contributes a GPU and is paid for
inference work,'' with the payment denominated in the network's native token and the fairness
enforced by the same consensus that finalizes the placement map. This is the point of contact
with the sovereign-cloud economics of the tenant-cloud companion~\cite{tenantcloud}, which
develops the wallet, staking, and settlement primitives this layer would sit on.

\section{Migration path: from LAN appliance to public network}
\label{sec:migration}

The design is only credible if there is a continuous path from what runs today to the network,
with no flag day. There is, in four stages, each shippable and independently valuable.

\paragraph{Stage 0 (today, \built{}).} Static \code{RING\_CONFIG} JSON, QR join token, raw-TCP
ring, canonical \code{f32} across vendors. Works on a trusted LAN; this is the baseline.

\paragraph{Stage 1 --- authenticate the wire (\design{}, smallest step).} Replace the ring's
\code{TcpStream} with a ZAP \code{WrapPQ} \code{net.Conn} and make the join token gate on a PQ
handshake (\S\ref{sec:threat}, Gaps 1--2). No consensus yet: the topology is still the static
map, but the wire is now authenticated and encrypted and the token is a credential. This alone
converts the appliance into something safe to run across untrusted networks, and it exercises the
engine's ZAP module end-to-end --- turning the \partialb{} \code{zap://} transport on
(\S\ref{sec:zap}) is precisely this step.

\paragraph{Stage 2 --- consensus owns the map (\design{}).} Stand up a Quasar chain whose state
is the membership and placement map. Nodes read topology from the chain instead of
\code{RingConfig::load()}; \code{hanzo cluster join} becomes an admission transaction rather than
a local launch. Placement is still operator-proposed and consensus-ratified (permissioned
membership, BFT config). The ring code is untouched; only the source of its topology moves from
file to chain.

\paragraph{Stage 3 --- permissionless miners and paid work (\design{}).} Open admission to
staked strangers, bind luminance to validated inference contribution, and turn on result
attestation and spot-check re-execution (\S\ref{sec:threat}--\S\ref{sec:econ}). This is the full
network: anyone bonds, joins, is placed by consensus, serves layers over authenticated ZAP, is
verified by a Photon committee, and is paid by luminance-weighted contribution.

The staging matters because each stage stands alone. Stage 1 is worth shipping even if Stages
2--3 never do (it fixes a real security gap). Stage 2 is worth shipping for a single operator's
multi-datacenter fleet (BFT config, no static files) without any economics. Only Stage 3 needs
the token and slashing machinery. There is no rewrite between stages --- the same ring, the same
ZAP, more of the control plane each time.

\section{Connection to the latency-physics companion}
\label{sec:latency}

A pipeline-parallel ring over a wide-area network has a latency budget a single box does not: each
pipeline stage adds a network hop to the critical path of every token. The companion
latency-physics analysis~\cite{latency} argues that this budget is spendable precisely because
two techniques compose over ZAP. First, \textbf{pipeline parallelism hides transport latency
under compute}: with enough micro-batches in flight, a stage's forward pass overlaps the next
stage's activation transfer, so the ring's throughput approaches the slowest stage's compute rate
rather than the sum of hop latencies --- and ZAP's zero-copy framing minimizes the per-hop
transport term to an AEAD seal, not a marshal. Second, \textbf{speculative decoding amortizes the
round-trip}: a draft model proposes several tokens that the ring verifies in one pipelined pass,
so the wide-area round-trip is paid once per \emph{accepted block} of tokens instead of once per
token. The engine already factors speculative decoding behind one \code{SpeculativeProposer}
trait~\cite{infstack}; over a consensus-managed ZAP ring, the draft can even run on a nearby,
brighter node while the target pipeline spans the fleet. The three-plane fabric is what makes
these composable: the data plane does the pipelining and speculation, the transport plane makes
each hop cheap and authenticated, and the control plane places draft and target where the latency
budget is best --- decomplected, so each optimization lives in exactly one plane. We cite the
companion for the quantitative budget; here it is enough that the architecture does not fight the
physics.

\section{Status: built, partial, and designed}
\label{sec:status}

Table~\ref{tab:status} is the honest ledger. It is the table to read first: it says exactly which
sentences of this paper are code on \code{main} and which are blueprint. The pattern is
deliberate --- every \emph{primitive} is built and cited; every \emph{composition} into the
network is designed. That is what makes this a buildable blueprint rather than a wish.

\begin{table}[t]
\centering
\small
\setlength{\tabcolsep}{5pt}
\begin{tabular}{@{}p{4.9cm}p{1.5cm}p{1.5cm}p{5.0cm}@{}}
\toprule
\textbf{Component} & \textbf{Plane} & \textbf{Status} & \textbf{Evidence / source} \\
\midrule
Cross-vendor inference ring (TP/PP, canonical \code{f32}) & Data & \built{} &
\code{distributed.rs}, \code{distributed/socket.rs}, \code{ops.rs}~\cite{infstack} \\
NCCL $+$ portable ring backends & Data & \built{} & \code{hanzo-quant/src/distributed/} \\
\code{hanzo cluster} scan-to-join (bs58/QR token) & Data & \built{} &
\code{hanzo-cli/src/cluster.rs} \\
Five-backend engine, decode parity, quant zoo & Data & \built{} &
companion stack paper~\cite{infstack} \\
ZAP Go wire (zero-copy, $17\times$/$345\times$ vs JSON-RPC) & Transport & \built{} &
\code{luxfi/zap}~\cite{zap,lp201} \\
ZAP PQ handshake, \code{WrapPQ}, verified \code{PeerID} & Transport & \built{} &
\code{conn\_pq.go} \\
ZAP engine module (in-process adapter surface) & Transport & \partialb{} &
\code{engine/src/zap/*.rs} (server/client) \\
\quad out-of-process \code{zap://} (capnp wire) & Transport & \partialb{} &
\code{engine/src/zap/transport.rs} (stub) \\
Quasar BFT engine; Photon, Wave, Focus & Control & \built{} &
\code{luxfi/consensus}~\cite{quasar,luxconsensus} \\
Luminance reputation (100 base, $\times1.1$/$\times0.9$) & Control & \built{} &
\code{photon/luminance.go} \\
BLS$+$Pulsar$+$ML-DSA threshold signing (QuasarCert) & Control & \built{} &
\code{protocol/quasar}~\cite{luxconsensus} \\
\midrule
Ring transport over ZAP (retire raw TCP) & Transport & \design{} & \S\ref{sec:transport}, Stage 1 \\
Join token as PQ handshake credential & Transport & \design{} & \S\ref{sec:threat}, Gap 2 \\
Consensus-owned membership \& placement map & Control & \design{} & \S\ref{sec:control}, Stage 2 \\
Permissionless miner admission (stake bond) & Control & \design{} & \S\ref{sec:econ}, Stage 3 \\
Luminance $\gets$ validated inference contribution & Control & \design{} & \S\ref{sec:control} \\
Result verification (threshold attest $+$ spot check) & Control & \design{} & \S\ref{sec:threat}, Gap 3 \\
Luminance-weighted reward / routing & Control & \design{} & \S\ref{sec:econ} \\
\bottomrule
\end{tabular}
\caption{Built-versus-designed ledger. Every primitive (upper block) is code on \code{main},
cited to its file with its measured cost where relevant; every network composition (lower block)
is design, cited to the section that specifies it and the migration stage that ships it.}
\label{tab:status}
\end{table}

\paragraph{Reproducing the built numbers.} The ZAP figures ($2.9$~ns parse, $17\times$/$345\times$
versus JSON-RPC, $114$K/$376$K TPS) reproduce with \code{go test -bench=. -benchmem} in
\code{luxfi/zap} and \code{luxfi/consensus/bench/}; the Photon/Wave/luminance costs ($3.03$~\textmu s,
$3.38$~\textmu s, $72$~ns) reproduce with \code{GOWORK=off go test -bench=. ./bench/} in
\code{luxfi/consensus}, measured on Apple M1 Max. The cross-vendor ring runs with a
\code{RING\_CONFIG} JSON per rank and \code{hanzo cluster init}/\code{join}; the cross-backend
engine numbers are reproduced by the protocol in the companion stack paper~\cite{infstack}. None
of these numbers is produced by this paper; each is attributed to its source so the blueprint's
foundation is auditable.

\section{Limitations}
\label{sec:limits}

The rigor is the persuasion, so the costs are as plain as the benefits.

\begin{itemize}[leftmargin=1.4em,itemsep=2pt]
\item \textbf{The network is a composition, not a shipped network.} The primitives are built and
cited; the assembly of \S\ref{sec:planes}--\S\ref{sec:econ} is \design{}. Nothing here should be
read as ``Hanzo runs a permissionless GPU network today.'' It runs a cross-vendor ring on a
trusted mesh (Stage 0) and ships every primitive the network needs.

\item \textbf{The ZAP engine wire is not yet turned on.} The Rust ZAP module is \partialb{}: its
in-process surface is exercised, but the \code{zap://} out-of-process transport is a documented
stub and a few methods return ``not yet wired'' (\S\ref{sec:zap}). Stage 1 depends on completing
it. The Go ZAP wire it mirrors is production-deployed, so the risk is integration, not invention.

\item \textbf{Wide-area tensor transport is unproven at scale.} The ring was built for a LAN.
Pipeline parallelism over a wide-area network is argued to be latency-tolerant
(\S\ref{sec:latency}) but not measured here; the companion latency paper carries that budget, and
a real deployment must validate it. Bandwidth, not just latency, may bound activation exchange
for large hidden dimensions.

\item \textbf{Result verification adds cost.} Threshold attestation and redundant re-execution
(\S\ref{sec:threat}) are not free: a QuasarCert verifies in $2$--$5$~ms and spot-check redundancy
duplicates work on sampled requests. The scheme is designed to make the \emph{expected} overhead
small (attest per epoch, spot-check a fraction) but it is a real tax the monolithic single-trust
model does not pay --- the price of permissionlessness.

\item \textbf{Confidentiality from the operator is out of scope.} This design authenticates and
verifies; it does not hide activations from the computing miner without a TEE tier
(\S\ref{sec:threat}). Sensitive workloads need that orthogonal extension.

\item \textbf{Luminance re-binding is a policy change with game-theoretic edges.} Moving the
\code{Illuminate} trigger from vote-success to inference-contribution is a one-line mechanism
change but a large incentive-design question: collusion among a committee, grinding on the
verification lottery, and reward-share disputes across a pipeline are real and not solved by the
primitive alone. We name them; we do not claim to have priced them.
\end{itemize}

\section{Conclusion}
\label{sec:conclusion}

The heterogeneous GPU ring already solves the hard, unglamorous problem --- one model computing
correctly across ROCm, CUDA, and Metal, bit-for-bit, over a canonical wire. What it lacks to
become a network is not more numerics but two planes it does not have: a control plane to decide
membership and placement among parties who do not trust each other, and a transport plane to
authenticate the strangers it talks to. Both already exist as best-in-class Lux primitives ---
Quasar consensus and the ZAP post-quantum wire --- and the entire contribution of this paper is
to hold the three concerns apart and compose them: consensus schedules, ZAP carries, the ring
computes. That decomplection is what a monolith cannot do and what makes ``any developer, any
miner, contributes a GPU and is paid for inference work'' a buildable system rather than a slogan.
We have drawn the line between built and designed down the middle of every claim: the ring, the
scan-to-join CLI, the ZAP wire with its verified \code{PeerID}, the Quasar engine with luminance
and threshold signing are code on \code{main}, cited and measured; the network that composes them
--- consensus-owned membership, the ZAP-transported ring, permissionless admission, and result
verification by threshold-signing over outputs --- is design, staged for a migration with no flag
day. The primitives are done. The composition is the blueprint. We publish it so the rest of the
AI-infrastructure community can build on the same three planes.

\begin{thebibliography}{20}
\small

\bibitem{infstack}
Hanzo AI --- ML Systems. \emph{Hanzo ML and Hanzo Engine: A Unified Multi-Backend, Multi-Modal
Inference Stack}. Hanzo AI white paper, 2026 (companion; the data-plane engine, cross-vendor ring,
and five-backend decode parity).

\bibitem{tenantcloud}
Hanzo AI Research. \emph{The Unified Sovereign-Tenant Cloud: Linear Shared-Nothing Scaling via
Per-Tenant SQLite and In-Process Composition}. Hanzo Industries white paper, 2026 (companion; the
\code{hanzo.network} fabric, wallet/staking, TEE trust tiers, and PQ identity primitives).

\bibitem{latency}
Hanzo AI --- ML Systems. \emph{The Latency Physics of Decentralized Decode: Pipeline Parallelism
and Speculative Decoding over ZAP}. Hanzo AI white paper, 2026 (companion; the wide-area latency
budget referenced in \S\ref{sec:latency}).

\bibitem{quasar}
Lux Industries. \emph{Quasar: Post-Quantum DAG-Native BFT Consensus} (LP-020 / LP-105). Lux
Network, 2026. \url{https://github.com/luxfi/consensus}.

\bibitem{luxconsensus}
Lux Industries. \emph{luxfi/consensus}: the Quasar engine --- \code{protocol/photon} (Fisher--Yates
committee selection, luminance), \code{protocol/wave} (FPC threshold voting), \code{protocol/focus},
\code{protocol/quasar} (BLS $+$ Pulsar $+$ ML-DSA threshold signing). Source and benchmarks
(\code{bench/}, \code{README.md}, \code{LLM.md}). \url{https://github.com/luxfi/consensus}.

\bibitem{zap}
Lux Industries. \emph{ZAP: Zero-copy Application Protocol} --- the universal wire.
\code{github.com/luxfi/zap} (\code{zap.go}, \code{conn\_pq.go} PQ handshake / \code{WrapPQ} /
\code{PeerID}, \code{README.md} benchmarks).

\bibitem{lp201}
Lux Industries. \emph{LP-201: ZAP as the Universal Wire} and \emph{LP-200: The ZAP Stack}. Lux
Proposals, 2026. \url{https://github.com/luxfi/lps} (the $17\times$/$345\times$ vs JSON-RPC and
$\sim$$4.5\times$ vs gRPC measurements).

\bibitem{mldsa}
National Institute of Standards and Technology. \emph{FIPS 204: Module-Lattice-Based Digital
Signature Standard (ML-DSA)}. NIST, 2024.

\bibitem{mlkem}
National Institute of Standards and Technology. \emph{FIPS 203: Module-Lattice-Based
Key-Encapsulation Mechanism Standard (ML-KEM)}. NIST, 2024.

\bibitem{bls}
D. Boneh, B. Lynn, and H. Shacham. Short Signatures from the Weil Pairing. \emph{Journal of
Cryptology}, 17(4):297--319, 2004.

\bibitem{petals}
A. Borzunov et al. \emph{Petals: Collaborative Inference and Fine-tuning of Large Models}. ACL
System Demonstrations, 2023 (decentralized inference over volunteer GPUs; no BFT control plane or
PQ-authenticated transport).

\bibitem{kwon2023}
W. Kwon et al. Efficient Memory Management for Large Language Model Serving with PagedAttention.
\emph{SOSP}, 2023 (the monolithic single-operator serving baseline).

\bibitem{pinkas}
B. Pinkas and M. Reiter, and the committee-grinding threat that motivates rejection-sampled
uniform selection in \code{photon/emitter.go}. (Bias-free sortition under a grinding adversary.)

\end{thebibliography}

\end{document}
